\documentclass[12pt]{article}

\usepackage{helpers}

\begin{document}

Name: \makebox[3in]{\hrulefill
%NAME
\hrulefill}

\vfill

\begin{center}
{\huge Test 3}

{\large 2026}
\end{center}

\testrules

\vfill

\pledge

\newpage

\section*{Instruction Pipelining}

Demonstrate understanding of instruction pipelining in the context of a CPU.

\begin{enumerate}

\item Without pipelining, how many clock cycles would it take for 3 instructions to execute? Explain your answer.
\vfill

\item With pipelining, how many clock cycles would it take for 3 instructions to execute? Explain your answer.
\vfill

\item What is a data hazard?
\vfill

\item What is operand forwarding?
\vfill

\newpage

\item What is a control hazard?
\vfill

\item What is eager execution?
\vfill

\item What is a pipeline stall?
\vfill

\end{enumerate}

\standardsfooter

\newpage

\section*{Stack Frames}

Demonstrate understanding of assembly code for stack frame setup and teardown. Be able to visualize and explain the code using memory diagrams. 

When drawing memory diagrams, begin at the top of \texttt{main} with \texttt{sp = 0x4004}, \texttt{fp = 0x4100}, and registers \texttt{r0} to \texttt{r3} empty. You do not need to track \texttt{pc} and \texttt{lr}. You may use placeholders such as ``old lr'' as needed.

The assembly code for these questions is found after the questions. 

\begin{enumerate}

\item Show the expected output of this program. You may choose any number of koalas you like.
\vfill

\item Draw a memory diagram showing the state of the program where it says \texttt{DIAGRAM 1 STOP HERE} below. Assume the user has 4 koalas. 
\vfill
\vfill
\vfill

\newpage


\item Draw a memory diagram showing the state of the program where it says \texttt{DIAGRAM 2 STOP HERE} below. Assume the user has 4 koalas. 
\vfill

\item Draw a memory diagram showing the state of the program where it says \texttt{DIAGRAM 3 STOP HERE} below. Assume the user has 4 koalas. 
\vfill

\end{enumerate}

\newpage 


\begin{verbatim}
@ global constants
.section .rodata
prompt: .ascii "How many koalas do you have? \0"
input_format: .ascii "%d\0"
reply: .ascii "That's %d thumbs!\n\0"

num_thumbs:
    @ stack frame setup, no local variables
    push {fp, lr}
    add fp, sp, #4
    @ ----------> DIAGRAM 2 STOP HERE
    @ multiply the number of koalas by 6, return result
    mov r2, #6
    mul r0, r0, r2
    @ stack frame teardown
    pop {fp, pc}

.text
.global main
main: 
    @ set up stack frame for main, one local variable
    push {fp, lr}
    add fp, sp, #4
    sub sp, sp, #4
    @ ask the user how many koalas they have
    ldr r0, prompt_ptr
    bl printf
    @ store response to local variable
    sub r1, fp, #8
    ldr r0, input_format_ptr
    bl scanf
    @ ----------> DIAGRAM 1 STOP HERE
    @ load local variable, pass it into num_thumbs
    ldr r0, [fp, #-8]
    bl num_thumbs    
    @ ----------> DIAGRAM 3 STOP HERE
    @ print the results
    mov r1, r0
    ldr r0, reply_ptr
    bl printf
    @ return 0, stack frame teardown
    mov r0, #0
    sub sp, fp, #4
    pop {fp, pc}

@ pointers
prompt_ptr: .word prompt
input_format_ptr: .word input_format
reply_ptr: .word reply
\end{verbatim}



\vfill

\standardsfooter

%END_STANDARD_HC

\end{document}